% GENERATED FILE. Edit canonical lessons, then run npm run generate:books.
% canonical-source-hash: fnv1a64:8555b2ecbd3ddd45
% canonical-lessons: TE-C479-pidugu, TE-C479-kundapota, TE-C479-kammu, TE-C479-masakabaru, TE-C479-tumpara

\chapter{Thunder and Downpours}
\label{ch:te-thunder-and-downpours}
\hlchaptermodality{479}

\begin{chapteropening}
\textbf{By the end of this chapter:} \emph{I can say a thunderbolt, a torrential rain, to close in, to grow dim and spray.}

Complete the last lesson of chapter 479: I can say a thunderbolt, a torrential rain, to close in, to grow dim and spray.
\end{chapteropening}

\section[\texorpdfstring{\te{పిడుగు} (piḍugu)}{పిడుగు (piḍugu)}]{\te{పిడుగు} (piḍugu) — a thunderbolt; a lightning strike}

\label{lesson:TE-C479-pidugu}



\begin{quote}
Before the new one: say the Telugu for a rainstorm, then the Telugu for to lose strength.
\end{quote}

\subsection*{You'll want to know: \te{పిడుగు}}
\textbf{\te{పిడుగు}} — \emph{piḍugu} — \textquotedblleft{}a thunderbolt; a lightning strike\textquotedblright{}.

\begin{quote}
\te{పొలంలో} \te{తాటిచెట్టు} \te{మీద} \te{పిడుగు} \te{పడింది}.
\end{quote}

\begin{quote}
\emph{polaṁlō tāṭiceṭṭu mīda piḍugu paḍindi.}
\end{quote}

\begin{quote}
Lightning struck a palm tree in the field.
\end{quote}

\te{పిడుగు} \te{పడు}, lightning strikes. Bad news that comes out of nowhere is also a \te{పిడుగు}.

\begin{cousinweb}
\textbf{\te{పిడుగు}} (\emph{piḍugu}) is native Dravidian: DEDR 4150 lists Telugu \emph{piḍugu}, \textquotedblleft{}a thunderbolt, lightning\textquotedblright{}, with Kannada \emph{piḍugu}.
\end{cousinweb}

\begin{grammarlens}[title={What you've built}]
One more word for getting around.
\end{grammarlens}

\subsection*{Your turn}
\begin{itemize}
\raggedright
  \item \emph{Say it:} \emph{piḍugu}
  \item \emph{Say it:} \emph{piḍugu}, once more
  \item \emph{Say it:} \emph{polaṁlō tāṭiceṭṭu mīda piḍugu paḍindi.} — Lightning struck a palm tree in the field.
  \item \emph{Recall:} say the Telugu for a rainstorm, then the Telugu for to lose strength, then say \emph{piḍugu} again
\end{itemize}

\begin{tcolorbox}[breakable,colback=teal!4,colframe=teal!35!black,title={Before you move on}]
What does \te{పిడుగు} mean? (\textquotedblleft{}A thunderbolt; a lightning strike\textquotedblright{}.) Say it once more.
\end{tcolorbox}
\section[\texorpdfstring{\te{కుండపోత} (kuṇḍapōta)}{కుండపోత (kuṇḍapōta)}]{\te{కుండపోత} (kuṇḍapōta) — a torrential rain}

\label{lesson:TE-C479-kundapota}



\begin{quote}
Before the new one: say the Telugu for to lose strength, then the Telugu for a thunderbolt.
\end{quote}

\subsection*{You'll want to know: \te{కుండపోత}}
\textbf{\te{కుండపోత}} — \emph{kuṇḍapōta} — \textquotedblleft{}a torrential rain\textquotedblright{}.

\begin{quote}
\te{సాయంత్రం} \te{కుండపోత} \te{వర్షం} \te{కురిసింది}.
\end{quote}

\begin{quote}
\emph{sāyantraṁ kuṇḍapōta varṣaṁ kurisindi.}
\end{quote}

\begin{quote}
In the evening it poured with rain.
\end{quote}

\te{కుండ}, a pot, + \te{పోత}, pouring: rain as if poured from a pot. Mostly \te{కుండపోత} \te{వర్షం}.

\begin{cousinweb}
\textbf{\te{కుండపోత}} (\emph{kuṇḍapōta}) is a Telugu compound built from \te{కుండ}, \textquotedblleft{}a pot\textquotedblright{}, and \te{పోత}, \textquotedblleft{}pouring\textquotedblright{}, the noun of \te{పోయు}, \textquotedblleft{}to pour\textquotedblright{} (DEDR 4407).
\end{cousinweb}

\begin{grammarlens}[title={What you've built}]
One more word for getting around.
\end{grammarlens}

\subsection*{Your turn}
\begin{itemize}
\raggedright
  \item \emph{Say it:} \emph{kuṇḍapōta}
  \item \emph{Say it:} \emph{kuṇḍapōta}, once more
  \item \emph{Say it:} \emph{sāyantraṁ kuṇḍapōta varṣaṁ kurisindi.} — In the evening it poured with rain.
  \item \emph{Recall:} say the Telugu for to lose strength, then the Telugu for a thunderbolt, then say \emph{kuṇḍapōta} again
\end{itemize}

\begin{tcolorbox}[breakable,colback=teal!4,colframe=teal!35!black,title={Before you move on}]
What does \te{కుండపోత} mean? (\textquotedblleft{}A torrential rain\textquotedblright{}.) Say it once more.
\end{tcolorbox}
\section[\texorpdfstring{\te{కమ్ము} (kammu)}{కమ్ము (kammu)}]{\te{కమ్ము} (kammu) — to close in; to envelop}

\label{lesson:TE-C479-kammu}



\begin{quote}
Before the new one: say the Telugu for a thunderbolt, then the Telugu for a torrential rain.
\end{quote}

\subsection*{You'll want to know: \te{కమ్ము}}
\textbf{\te{కమ్ము}} — \emph{kammu} — \textquotedblleft{}to close in; to envelop\textquotedblright{}.

\begin{quote}
\te{ఆకాశంలో} \te{నల్లని} \te{మబ్బులు} \te{కమ్మాయి}.
\end{quote}

\begin{quote}
\emph{ākāśaṁlō nallani mabbulu kammāyi.}
\end{quote}

\begin{quote}
Dark clouds closed in across the sky.
\end{quote}

Past (it …): \te{కమ్మింది}, \emph{kammindi}.

Said of clouds, smoke or darkness that cover everything. Its subject is a thing, so the past shown is the third person: \te{కమ్మింది}, it closed in.

\begin{cousinweb}
\textbf{\te{కమ్ము}} (\emph{kammu}) is native Dravidian: DEDR 1221 lists Telugu \emph{kammu}, \textquotedblleft{}to cover, to surround, to overspread\textquotedblright{}, beside \emph{kappu}, \textquotedblleft{}to spread over in a thick covering, as clouds\textquotedblright{}.
\end{cousinweb}

\begin{grammarlens}[title={What you've built}]
One more word for getting around.
\end{grammarlens}

\subsection*{Your turn}
\begin{itemize}
\raggedright
  \item \emph{Say it:} \emph{kammu}
  \item \emph{Say it:} \emph{kammu}, once more
  \item \emph{Say it:} \emph{ākāśaṁlō nallani mabbulu kammāyi.} — Dark clouds closed in across the sky.
  \item \emph{Recall:} say the Telugu for a thunderbolt, then the Telugu for a torrential rain, then say \emph{kammu} again
\end{itemize}

\begin{tcolorbox}[breakable,colback=teal!4,colframe=teal!35!black,title={Before you move on}]
What does \te{కమ్ము} mean? (\textquotedblleft{}To close in; to envelop\textquotedblright{}.) Say it once more.
\end{tcolorbox}
\section[\texorpdfstring{\te{మసకబారు} (masakabāru)}{మసకబారు (masakabāru)}]{\te{మసకబారు} (masakabāru) — to grow dim; to blur}

\label{lesson:TE-C479-masakabaru}



\begin{quote}
Before the new one: say the Telugu for a torrential rain, then the Telugu for to close in.
\end{quote}

\subsection*{You'll want to know: \te{మసకబారు}}
\textbf{\te{మసకబారు}} — \emph{masakabāru} — \textquotedblleft{}to grow dim; to blur\textquotedblright{}.

\begin{quote}
\te{వర్షంలో} \te{దూరం} \te{అంతా} \te{మసకబారింది}.
\end{quote}

\begin{quote}
\emph{varṣaṁlō dūraṁ antā masakabārindi.}
\end{quote}

\begin{quote}
In the rain everything in the distance grew dim.
\end{quote}

Past (it …): \te{మసకబారింది}, \emph{masakabārindi}.

\te{మసక}, hazy, + \te{బారు}, to become. Old glasses and old eyes also \te{మసకబారు}. Its subject is a thing, so the past shown is the third person: \te{మసకబారింది}, it grew dim.

\begin{cousinweb}
\textbf{\te{మసకబారు}} (\emph{masakabāru}) is a Telugu compound: \te{మసక}, \textquotedblleft{}dim, hazy\textquotedblright{}, native Dravidian (DEDR 4627), and -\te{బారు}, \textquotedblleft{}to become\textquotedblright{}, whose origin is not settled.
\end{cousinweb}

\begin{grammarlens}[title={What you've built}]
One more word for getting around.
\end{grammarlens}

\subsection*{Your turn}
\begin{itemize}
\raggedright
  \item \emph{Say it:} \emph{masakabāru}
  \item \emph{Say it:} \emph{masakabāru}, once more
  \item \emph{Say it:} \emph{varṣaṁlō dūraṁ antā masakabārindi.} — In the rain everything in the distance grew dim.
  \item \emph{Recall:} say the Telugu for a torrential rain, then the Telugu for to close in, then say \emph{masakabāru} again
\end{itemize}

\begin{tcolorbox}[breakable,colback=teal!4,colframe=teal!35!black,title={Before you move on}]
What does \te{మసకబారు} mean? (\textquotedblleft{}To grow dim; to blur\textquotedblright{}.) Say it once more.
\end{tcolorbox}
\section[\texorpdfstring{\te{తుంపర} (tumpara)}{తుంపర (tumpara)}]{\te{తుంపర} (tumpara) — spray; fine drizzle}

\label{lesson:TE-C479-tumpara}



\begin{quote}
Before the new one: say the Telugu for to close in, then the Telugu for to grow dim.
\end{quote}

\subsection*{You'll want to know: \te{తుంపర}}
\textbf{\te{తుంపర}} — \emph{tumpara} — \textquotedblleft{}spray; fine drizzle\textquotedblright{}.

\begin{quote}
\te{జలపాతం} \te{దగ్గర} \te{తుంపర} \te{మా} \te{మీద} \te{పడింది}.
\end{quote}

\begin{quote}
\emph{jalapātaṁ daggara tumpara mā mīda paḍindi.}
\end{quote}

\begin{quote}
Near the waterfall the spray fell on us.
\end{quote}

Fine drops of water: from a waterfall, a wave or a light rain.

\begin{cousinweb}
\textbf{\te{తుంపర}} (\emph{tumpara}) is native Dravidian: DEDR 3398 lists Telugu \emph{tuppara}, \textquotedblleft{}a drop of water, a thin drizzle\textquotedblright{}, with the variant \emph{tumpara}.
\end{cousinweb}

\begin{grammarlens}[title={What you've built}]
One more word for getting around.
\end{grammarlens}

\subsection*{Your turn}
\begin{itemize}
\raggedright
  \item \emph{Say it:} \emph{tumpara}
  \item \emph{Say it:} \emph{tumpara}, once more
  \item \emph{Say it:} \emph{jalapātaṁ daggara tumpara mā mīda paḍindi.} — Near the waterfall the spray fell on us.
  \item \emph{Recall:} say the Telugu for to close in, then the Telugu for to grow dim, then say \emph{tumpara} again
\end{itemize}

\begin{tcolorbox}[breakable,colback=teal!4,colframe=teal!35!black,title={Before you move on}]
What does \te{తుంపర} mean? (\textquotedblleft{}Spray; fine drizzle\textquotedblright{}.) Say it once more.
\end{tcolorbox}
